\documentclass[12pt]{article}
\usepackage[bookmarks=false]{hyperref}
\usepackage{enumerate}
\usepackage{amssymb,amsmath,amsthm}
\usepackage{graphicx}
\usepackage{geometry}

%\pagestyle{empty}

\begin{document}

\section*{Hamiltonian Alternatives}

The working Hamiltonian is given by
\[
H = -\frac{k}{N}\sum_{i\neq j}J_{ij}\sigma_i\sigma_j,
\]
where $J_{ij}$ is the interaction strength between spins $i$ and $j$ with each spin associated with one of $k$ targets (so there are $k$ groups of spins) and existing in one of two states: $\sigma_i=0$ or $\sigma_i=1$ - e.g., a neuron that is not firing or firing. The sum over all $i\neq j$ assumes a fully-connected network of spins/neurons (all interact with all), and only pairs of ``on'' spins contribute to the energy. A positive $J_{ij}$ indicates an excitatory interaction between spin $i$ and spin $j$ which lowers energy (consensus lowers the energy of the system) while a negative $J_{ij}$ indicates an inhibitory interaction that raises the energy.

Note: in Grobonos's PRXLife paper, he also multiplies by a constant representing velocity. This would imply that the faster an organism is moving, the steeper the energy landscape and again, the less ability to explore the space. Potentially quite biological! 

If we consider instead spins that can be -1, 0, or 1, but everything else is the same, the problem is that two -1 spins interact to lower the energy, and the -1 spins are interacting with the targets. Not good.

Instead, what if we have something like this:
\[
H = -\frac{k_t}{N_t}\sum_{i\neq j}J^t_{ij}\sigma_{t,i}\sigma_{t,j} + \frac{k_b}{N_b}\sum_{i\neq j}J^b_{ij}\sigma_{b,i}\sigma_{b,j}
\]
where spins $t$ are associated with one of $k_t$ targets while spins $b$ are associated with one of $k_b$ blocking obstacles. All spins are either 1 or 0, again corresponding to firing or not firing respectively.

So, what is function of the interaction kernel? I think it should be a Gaussian, where the width of the Gaussian is based on the width of the organism navigating the landscape.

Distance becomes really important here. Obstacles that are far away aren't obstacles, and they become more of a priority the closer they get. This may be a good lens through which to examine targets as well, but targets remain targets at any difference. There's just a question of perception - if they are far enough away, they may not be identified in a collective behavior scenario.

\end{document}